\documentclass[10pt,a4paper]{amsart}
\usepackage[margin=19mm]{geometry}
\usepackage{amsmath,amssymb,mathtools}
\usepackage[scr=rsfs,cal=euler]{mathalfa}
\usepackage{xfrac}
\usepackage[unicode,hidelinks,bookmarks=false]{hyperref}
\newcommand{\ie}{i.e.\ }
\newcommand{\cf}{cf.\ }
\newcommand{\resp}{respectively\ }
\newcommand{\Subsection}{\subsection}
\providecommand{\nobreakdash}{\nobreak\hbox{-}}
\setlength{\emergencystretch}{2em}
\multlinegap0pt
\usepackage{bm}
\usepackage{bbm}
\usepackage{dsfont}
\usepackage{xspace}
\usepackage{cancel}
\usepackage{mathtools}
\usepackage{amsthm}
\makeatletter
\@ifundefined{theorem}{\newtheorem{theorem}{Theorem}}{}
\@ifundefined{lemma}{\newtheorem{lemma}[theorem]{Lemma}}{}
\@ifundefined{proposition}{\newtheorem{proposition}[theorem]{Proposition}}{}
\makeatother
\newcommand{\R}{\mathbb{R}}
\newcommand{\Z}{{\mathbb{Z}}}
\title[Global dynamics above the ground state for the energy-critic]{Global dynamics above the ground state for the energy-critical Hartree equation with radial data}
\author{Xuemei Li, Chenxi Liu, Guixiang Xu}
\date{Source: arXiv:2506.05406v1 (CC BY 4.0)}
\begin{document}
\maketitle

\noindent\emph{Source and licence.} Global dynamics above the ground state for the energy-critical Hartree equation with radial data. Version arXiv:2506.05406v1 (CC BY 4.0), distributed under the Creative Commons Attribution 4.0 International licence, as recorded in the arXiv licence metadata for this exact version and shown on the abstract page https://arxiv.org/abs/2506.05406v1. Full rights notice: licenses/arxiv-2506.05406-CC-BY-4.0.txt.\par\smallskip

\noindent\textit{Editorial scope.} Counted unit: the Proposition whose source label is \texttt{prop:nlr dist} in the article (source0.tex lines 897--909, source label \texttt{prop:nlr dist}), delivered together with the article's own complete original proof of it (source0.tex lines 911--954, one \texttt{proof} environment of 44 lines). No mathematics is written, repaired or re-derived: statement and proof are the article's own text line for line. Required context retained uncounted: equ:Hart (lines 245--248); quan:energy (lines 274--277); quan:KIG (lines 545--551); est:EK expa (lines 600--611); prop:orth decom (lines 644--660); prop:spect decom (lines 713--732); decom:v+ (lines 714--716); prop:orth cntl (lines 824--828); quan:equi nor (lines 834--836); prop:unif lwp (lines 855--859); quan:d\_0 (lines 878--880). Every definition, display and label of those lines is retained unchanged. Omitted from this package and not redistributed: 1--244, 249--273, 278--544, 552--599, 612--643, 661--712, 717--823, 829--833, 837--854, 860--877, 881--896, 910--910, 955--2682. Nothing inside a retained excerpt is omitted or altered: every retained body is delivered exactly as the article's lines are written, so no omission receipt of any kind accompanies this package. Citations: the retained excerpts carry no citation command at all, so no bibliography entry of the article is transcribed and no reference is redistributed. Reference adaptation: none was needed and none was made; every reference inside the delivered unit resolves inside it, so no unresolved reference or citation is produced. Printed slips kept literal: no printed word, symbol, hypothesis or number of the counted statement or of its proof was altered. Layout and class: the article's own class is \texttt{amsart} with packages including amssymb, amsmath, mathtools, ifpdf, vmargin, mathrsfs, epstopdf, setspace; the wrapper is the lane's pinned \texttt{amsart} editorial wrapper at 10pt on A4, which restates the theorem-like environments the retained text needs and adds the article's own macro definitions that the retained text needs (\texttt{\textbackslash R}, \texttt{\textbackslash Z}), with extra wrapper packages (bm, bbm, dsfont, xspace, cancel, mathtools). The delivered numbering is the unit's own. Assets: no retained span contains a figure environment, a graphics-inclusion command, a PDF-page inclusion, a table or a TikZ picture; no image file is copied into this package, the package declares no asset dependency and no image file is redistributed with it. Licence: version arXiv:2506.05406v1 of this article is distributed under the Creative Commons Attribution 4.0 International licence, as recorded in the arXiv licence metadata for this exact version and shown on the abstract page https://arxiv.org/abs/2506.05406v1. A full rights notice is delivered at licenses/arxiv-2506.05406-CC-BY-4.0.txt. Characters: every retained excerpt is pure ASCII, so the delivered engine, XeLaTeX with its default Latin Modern text font, reports no missing character.
	\begin{equation}\label{equ:Hart}
	i \partial_t u - \Delta u  - \left( |x|^{-4}*|u|^2\right)u =0,\quad
	(t,x)\in \R \times \R^d,
	\end{equation}

	\begin{align}\label{quan:energy}
 E(u(t)):=&\frac12 \int_{\R^d} \big| \nabla u(t,x)\big|^2 dx-\frac{1}{4} \iint_{\R^d\times \R^d} \frac{\big|u(t,x)\big|^2\big|u(t,y)\big|^2}{|x-y|^4} dxdy\\
	=&E(u_0) \nonumber
	\end{align}

		\begin{equation}\label{quan:KIG}
		\begin{aligned}
		K(\varphi):=&\Vert \nabla \varphi \Vert_{L^2}^{2} -{\Vert{\left(|x|^{-4}*|\varphi|^2\right)|\varphi|^2}\Vert}_{L^1},\\
		G(\varphi):=&E(\varphi)-K(\varphi)/4=\Vert \nabla \varphi \Vert_{L^2}^{2} /4,\\
		I(\varphi):=&E(\varphi)-K(\varphi)/2={\Vert{\left(|x|^{-4}*|\varphi|^2\right)|\varphi|^2}\Vert}_{L^1}/4.
		\end{aligned}
		\end{equation}

	\begin{lemma} For the functionals defined in \eqref{quan:energy} and \eqref{quan:KIG}, we have
		\begin{equation}\label{est:EK expa}
		\begin{aligned}
		E(W+v)=&E(W)+\frac{1}{2}\langle \mathcal{L}v,v\rangle-C(v),\\
		K(W+v)=&-2\langle \nabla W,\nabla v\rangle+O(\Vert v\Vert_{\dot H^1}^2),
		\end{aligned}
		\end{equation}
		where 
		\begin{equation*}
		C(v):=\frac{1}{4}\iint \frac{\big|v(t,x)\big|^2\big|v(t,y)\big|^2+4W(x)v_1(t,x)\big|v(t,y)\big|^2}{|x-y|^4} dxdy=O(\Vert v \Vert_{\dot H^1}^3).
		\end{equation*}
	\end{lemma}

	\begin{proposition}(Orthogonal decomposition of $\varphi$)\label{prop:orth decom}. There exist an absolute constant $0<\delta_E\ll1$ and a $C^1$ function $\left(\widetilde{\theta},\widetilde{\sigma} \right):B_{\delta_E}(\mathcal{W}) \rightarrow (\R/2\pi\Z)\times \R$ with the following properties. For any $\varphi\in B_{\delta_E}(\mathcal{W})$, putting 
		\begin{equation}\label{decom:var}
		\varphi=e^{i\widetilde{\theta}(\varphi)}S_{-1}^{\widetilde{\sigma}(\varphi)}(W+v),
		\end{equation}	
		we have 
		\begin{equation*}
		\begin{aligned}
		(v,iW)_{\dot H^1}=(v,\widetilde{W})_{\dot H^1}=0, \; d_{\mathcal{W}}(\varphi)\sim \Vert v\Vert_{\dot H^1}.
		\end{aligned}
		\end{equation*}
		Moreover, $\left(\widetilde{\theta}(\varphi),\widetilde{\sigma} (\varphi)\right) \in  (\R/2\pi\Z)\times \R$ is unique for the above property. 
		
		Furthermore, if $\Vert \varphi-W_{\theta,\sigma}\Vert_{\dot H^1} \ll 1$ for some $(\theta,\sigma)\in\R^2$, then 
		\begin{equation*}
		\big|\left(e^{i\widetilde{\theta}(\varphi)}-e^{i\theta}, \widetilde{\sigma}(\varphi)-\sigma \right)\big|\lesssim \Vert \varphi-W_{\theta,\sigma}\Vert_{\dot H^1}.
		\end{equation*}
	\end{proposition}

	\begin{proposition}(Spectral decomposition of $v$)\label{prop:spect decom}. For any $v \in \dot H^1$, there exists a unique decomposition
		\begin{equation}\label{decom:v+}
		v=\lambda_+g_{+} + \lambda_-g_{-}+\gamma,\;\;\lambda_{\pm}\in\R,\;\;\gamma\in  \dot H^1,
		\end{equation}
		such that $\omega(g_{\pm},\gamma)=0$ (This implies that $\gamma\in G_{\bot}$).
		
		After normalizing $g_{\pm}$ (or $g_1$ and $g_2$) such that 
		\begin{equation}\label{con:ome}
		\omega(g_{+},g_{-})=2\langle g_{1},g_{2} \rangle=1,\; \pm\omega(W,g_{\mp})=\langle W,g_2 \rangle>0,
		\end{equation}
		the above decomposition is given by
		\begin{equation*}
		\lambda_{\pm}=\pm \omega(v,g_{\mp}).
		\end{equation*}
		Putting $\lambda_1:=(\lambda_+ +\lambda_-)/2$ and $\lambda_2:=(\lambda_+ -\lambda_-)/2$, it can also be written as 
		\begin{equation}\label{decom:v1}
		v=2\lambda_{1}g_1-2i\lambda_{2}g_2+\gamma,\; \lambda_1=(v_1|g_2),\; \lambda_2=-(v_2|g_1),
		\end{equation}
		with $(\gamma_1|g_2)=(\gamma_2|g_1)=0$.
	\end{proposition}

	\begin{proposition}(Control of orthogonal direction)\label{prop:orth cntl}. For function $\gamma$ defined in \eqref{decom:v+}, we have
		\begin{equation*}
		\Vert \gamma \Vert_{\dot H^1}^2 \sim \langle \mathcal{L}\gamma,\gamma \rangle.
		\end{equation*}
	\end{proposition}

	\begin{equation}\label{quan:equi nor}
	\Vert v \Vert_{E}^2:=\mu (\lambda_{1}^{2}+\lambda_{2}^{2})+\frac{1}{2}\langle \mathcal{L}\gamma,\gamma \rangle \sim \lambda_{1}^{2}+\lambda_{2}^{2}+\Vert \gamma \Vert_{\dot H^1}^2 \sim \Vert v \Vert_{\dot H^1}^2.
	\end{equation}

	\begin{proposition}(Uniform local existence in $\tau$)\label{prop:unif lwp}. There exists an absolute constant $\delta_L\in(0,\delta_{E}/2)$ such that for any solution $u$ of \eqref{equ:Hart} with $d_{\mathcal{W}}(u(0))=:\delta\in\left[ 0,2\delta_L \right]$, we have $T_{\pm}(u)>3e^{-2\sigma(0)}=:T_0$, $\pm\left( \tau(\pm T_0)-\tau(0) \right)>2$, and for $|t|\leq T_0$,
		\begin{equation*}
		\delta_{E}>d_{\mathcal{W}}(u(t))\sim\delta,\;\; \sigma(t)=\sigma(0)+O(\delta).
		\end{equation*}
	\end{proposition}

	\begin{equation}\label{quan:d_0}
	d_0(\varphi)^2:=E(\varphi)-E(W)+2\mu\lambda_{1}^2.
	\end{equation}

	\begin{proposition}(Nonlinear distance function)\label{prop:nlr dist}. The function $\widetilde{d}_{\mathcal{W}}$ on $\dot H^1_{rad}$ is invariant for the rotation and scaling, and equivalent to $d_{\mathcal{W}}$. Precisely, there exists an absolute constant $C\in\left(1,\infty\right)$ such that for all $\varphi\in \dot H^1_{rad}$ and $(\alpha,\beta)\in \R^2$,
		\begin{equation*}
		d_{\mathcal{W}}(\varphi)/C \leq \widetilde{d}_{\mathcal{W}}(\varphi)=\widetilde{d}_{\mathcal{W}}(e^{i\alpha}S_{-1}^{\beta}\varphi) \leq C	d_{\mathcal{W}}(\varphi).
		\end{equation*}
Moreover, there  exists an absolute constant $c_D\in\left(0,1\right)$ such that putting
		\begin{equation*}
		\check{\mathcal{H}}:=\{ \varphi\in \dot H^1_{rad}| E(\varphi)<E(W)+(c_D\widetilde{d}_{\mathcal{W}}(\varphi))^2\},
		\end{equation*}
		we have
		\begin{equation}\label{est:dist eigen}
		\varphi\in B_{\delta_L}(\mathcal{W}) \cap \check{\mathcal{H}} \Longrightarrow \widetilde{d}_{\mathcal{W}}(\varphi) \sim |\lambda_{1}|.
		\end{equation}
	\end{proposition}

	\begin{proof}
		First we proof $\widetilde{d}_{\mathcal{W}} \sim d_{\mathcal{W}}$. If $d_{\mathcal{W}}(\varphi)\geq 2\delta_L$, then $\widetilde{d}_{\mathcal{W}}=d_{\mathcal{W}}$ by the defination of $\widetilde{d}_{\mathcal{W}}(\varphi)$. It remains to consider the case $d_{\mathcal{W}}(\varphi)\leq 2\delta_L$. Since $\varphi\in B_{2\delta_L}(\mathcal{W})\subset B_{\delta_E}(\mathcal{W})$, we can  decompose $\varphi$ by Proposition \ref{prop:orth decom} and \ref{prop:spect decom}. By \eqref{est:EK expa} and the scale invariance of $E$, we have
		\begin{equation}\label{est:E expa}
		E(\varphi)-E(W)=E(W+v)-E(W)=\frac{1}{2}\langle \mathcal{L}v,v \rangle-C(v)=-\mu\lambda_+\lambda_-+\frac{1}{2}\langle \mathcal{L}\gamma,\gamma \rangle-C(v).
		\end{equation}
		Hence, using Proposition \ref{prop:orth cntl} and \eqref{quan:d_0},
		\begin{equation*}
		d_0(\varphi)^2=\Vert v \Vert_{E}^2-C(v)=\Vert v \Vert_{E}^2+O(\Vert v \Vert_{\dot H^1}^3)\sim d_{\mathcal{W}}(\varphi)^2.
		\end{equation*}
		Then by  Proposition \ref{prop:unif lwp}, we have $d_1(\varphi)\sim d_{\mathcal{W}}(\varphi)$, and so $\widetilde{d}_{\mathcal{W}}(\varphi)\sim d_{\mathcal{W}}(\varphi)$.
		
		Next we prove \eqref{est:dist eigen}. On the one hand, since $d_{\mathcal{W}}(\varphi)<\delta_{L}$ and $E(\varphi)-E(W)<(c_D\widetilde{d}_{\mathcal{W}}(\varphi))^2$, we have
		\begin{equation*}
		\widetilde{d}_{\mathcal{W}}(\varphi)^2 \sim d_0(\varphi)^2=E(\varphi)-E(W)+2\mu\lambda_{1}^2,
		\end{equation*}
		and so $\widetilde{d}_{\mathcal{W}}(\varphi)^2 \lesssim \lambda_{1}^2$. On the other hand, we see that $\widetilde{d}_{\mathcal{W}}(\varphi)^2\sim \Vert v \Vert_{E}^2 \gtrsim \lambda_{1}^2$ by \eqref{quan:equi nor}.
		
		Finally, we check the invariance for the rotation and scaling. Let $(\alpha,\beta)\in\R^2$, $\varphi\in B_{\delta_E}(\mathcal{W})$ and let $u$ and $u'$ be the solutions of \eqref{equ:Hart} with the initial data
		\begin{equation*}
		u(0)=\varphi,\;\; u'(0)=e^{i\alpha} S_{-1}^{\beta}\varphi,
		\end{equation*}
		with the decompositions by Proposition \ref{prop:orth decom} and the rescaled time functions
		\begin{align*}
		u=&e^{i\theta} S_{-1}^{\sigma}(W+v),\;\;u'=e^{i\theta'} S_{-1}^{\sigma'}(W+v'),\\ \dot\tau=&e^{2\sigma},\;\;\dot\tau'=e^{2\sigma'},\;\;\tau(0)=0=\tau'(0).
		\end{align*}
		Then 
		\begin{equation*}
		e^{i\theta'(0)} S_{-1}^{\sigma'(0)}(W+v'(0))=e^{i\alpha} S_{-1}^{\beta}\varphi=e^{i\alpha}e^{i\theta(0)}S_{-1}^{\beta} S_{-1}^{\sigma(0)}(W+v(0)),
		\end{equation*}
		that is 
		\begin{equation*}
		\theta'(0)=\theta(0)+\alpha,\;\;\sigma'(0)=\sigma(0)+\beta.
		\end{equation*}
		By the uniqueness of $\left(\widetilde{\theta},\widetilde{\sigma} \right)$ in Proposition \ref{prop:orth decom}, we have $\left(\theta',\sigma' \right)=\left(\theta,\sigma\right)+(\alpha,\beta)$.
		Integrating the following identity
		\begin{equation*}
		\dot\tau'=e^{2\sigma'}=e^{2\sigma+2\beta}=e^{2\beta}\dot\tau,
		\end{equation*}
		we have $\tau'=\tau e^{2\beta}$, while the invariance of the equation \eqref{equ:Hart} implies $u'(t)=e^{i\alpha} S_{-1}^{\beta}u(e^{2\beta}t)$. Hence $v$ is invariant in the rescaled time, namely
		\begin{equation*}
		\tau(t)=\tau'(t') \Longrightarrow v(t)=v'(t'),
		\end{equation*}
		which is inherited by $\lambda_{\pm},\lambda_1,\lambda_2$ and $\gamma$. Therefore $d_0$ and $d_1$ are invariant, so is $\widetilde{d}_{\mathcal{W}}$.
	\end{proof}

\end{document}
